\chapter{Background and Related Work}

\section{Termination Analysis Techniques}

The theory of termination analysis relies on the notion of well-founded
relations. A binary relation $\prec$ on a set $S$ is well-founded if there are
no infinite descending chains. That is, there exists no infinite sequence
$(s_n)_{n\in\mathbb{N}}$ of elements in $S$ such that $s_{n+1} \prec s_n$
for all $n\in\mathbb{N}$ \cite[31]{programSemantics}.
If every recursive call can be shown to pass arguments that are strictly
smaller than the parameters for some well-founded relation, the function
terminates.
The following techniques, commonly used in existing tools, extend this idea.

\subsection{Lexicographic Ordering}
In functions with multiple arguments, we can construct a well-founded relation
for the tuple of arguments from well-founded relations on each argument. 
Formally, given well-founded relations $\prec_i$ on sets $A_i$, the
lexicographic order $\prec_\text{lex}$ on $A_1\times\cdots\times A_n$ is defined
as \cite[175]{programSemantics}:
$$ (a_1, \ldots, a_n) \prec_\text{lex} (b_1, \ldots, b_n) \iff
\exists i\in\{1, \ldots, n\} \text{ such that } a_i \prec_i b_i
\text{ and } \forall j<i, a_j = b_j .$$


\subsection{Size-Change Principle}
The size-change principle \cite{SizeChange} is a global termination analysis
techniques that models how the size of values change at each function call.
The analysis abstracts a program into a finite set of
\emph{size-change graphs}, where each graph represents a function call and
captures the size relations between the caller's and callee's parameters.
Size change termination is specified as a decision problem: given a finite set
of size-change graphs, does every infinite sequence of calls contain an infinite
decreasing thread of parameters? If the answer is yes, the program must
terminate, as well-founded data cannot decrease indefinitely.
Although quite effective, the algorithm to solve this decision problem requires
computing the transitive closure of the graphs which is computationally
expensive and has been shown to be \textsf{PSPACE-complete}.

\subsection{User-Provided Proofs}
When automatic analysis fails, systems fall back on human intervention. In proof
assistant environment, these are mathematical proofs. In the context of this
work, this technique takes the form of escape hatches.

\section{Existing Tools}
The challenge of verifying program termination has led to the development of
various tools, ranging from proof assistants to automated analyzers.

\subsection{Rocq, Agda, and Lean}
Rocq \cite{rocq}, Agda \cite{agda}, and Lean \cite{lean}
are state-of-the-art interactive proof assistants based on dependent type
theory. They enforce termination for all functions by default.
Rocq uses a syntactic guard condition for fixpoints, Agda and Lean employ a more
sophisticated termination checker based on the above techniques.
They all fall back on manually provided proofs, which is often too
labor-intensive for general-purpose software engineering.

\subsection{AProVE}
AProVE \cite{AProVE} is an automated system for proving termination for
general-purpose languages such as Java, C, and Haskell. It has achieved
state-of-the-art results in the Termination Competition
\footnote{\url{https://termination-portal.org/wiki/Termination_Competition_2025}}.
For Haskell, \cite{AProVEHaskell} implements in AProVE an approach, that
symbolically evaluates a given start term under Haskell's lazy evaluation to
construct a termination graph, then translates that graph into dependency-pair
problems solved by term rewriting techniques.

\subsection{Stainless}
Stainless \cite{DBLP:journals/corr/abs-1904-03482} is a verification framework
specifically designed for Scala. It translates Scala code into logical
statements, verified by an underlying SMT solver. Stainless can automatically
prove intricate mathematical bounds on integers, infer well-founded relations to
use, and has semantic awareness of values. However, Stainless operates as a
separate toolchain, requires external dependencies and adds significant
verification time.

